% !TeX spellcheck = es
\documentclass[preprint,3p]{elsarticle}

\usepackage{lineno,hyperref}
\usepackage{tabularx}
\usepackage{multirow}
\usepackage{listings}
\usepackage{color,soul}

%\usepackage{subfigure}

\usepackage{graphicx}
\usepackage{subcaption}

\modulolinenumbers[5]

\journal{Future Generation Computer Systems}

%\newcommand\boldred[1]{\textcolor{red}{\textbf{#1}}}
%\newcommand\boldblue[1]{\textcolor{blue}{\textbf{#1}}}
%\newcommand\boldcyan[1]{\textcolor{cyan}{\textbf{#1}}}

\newcommand\boldred[1]{#1}
\newcommand\boldblue[1]{#1}
\newcommand\boldcyan[1]{#1}


%%%%%%%%%%%%%%%%%%%%%%%
%% Elsevier bibliography styles
%%%%%%%%%%%%%%%%%%%%%%%
%% To change the style, put a % in front of the second line of the current style and
%% remove the % from the second line of the style you would like to use.
%%%%%%%%%%%%%%%%%%%%%%%

%% Numbered
%\bibliographystyle{model1-num-names}

%% Numbered without titles
%\bibliographystyle{model1a-num-names}

%% Harvard
%\bibliographystyle{model2-names.bst}\biboptions{authoryear}

%% Vancouver numbered
%\usepackage{numcompress}\bibliographystyle{model3-num-names}

%% Vancouver name/year
%\usepackage{numcompress}\bibliographystyle{model4-names}\biboptions{authoryear}

%% APA style
%\bibliographystyle{model5-names}\biboptions{authoryear}

%% AMA style
%\usepackage{numcompress}\bibliographystyle{model6-num-names}

%% `Elsevier LaTeX' style
\bibliographystyle{elsarticle-num}
%%%%%%%%%%%%%%%%%%%%%%%


\begin{document}

\begin{frontmatter}

\title{
	Blockchain-based resource accounting for a serverless execution platform
}

%% Group authors per affiliation:
\author{Jonatan Enes\corref{cor}}
\ead{jonatan.enes@udc.es}
\author{Roberto R. Exp\'osito\corref{cor2}}
\ead{rreye@udc.es}
\author{Juan Touri\~no\corref{cor2}}
\ead{juan@udc.es}

%[gac]
\address{Universidade da Coru\~na, CITIC, Computer Architecture Group, Campus de A Coru\~na, Spain}

\cortext[cor]{Corresponding author}

\begin{abstract}

\end{abstract}

\begin{keyword}
anonymous computing, decentralized accounting
\end{keyword}

\end{frontmatter}

\linenumbers

\section{Introduction}
\label{sec:intro}

\section{Experiments}
\label{sec:experiments}

\section{Común a todos los experimentos}

\begin{itemize}
	\item Se usaría 1 solo contenedor, dentro de 1 app y asociada a 1 usuario. El objetivo del artículo no sería demostrar tanto la plataforma serverless y sus capacidades de gestión de múltiples contenedores (eso ya lo han hecho otros artículos), sino sus funciones de contabilidad e integración con la blockchain.
	
	\item Se definirá 1 o 2 aplicaciones y se usarán en 3 escenarios de gestión del crédito distintos.
\end{itemize}

\newcommand{\ContainerUp}{Se levanta un contenedor sin nada ejecutándose, con unos recursos asignados iniciales}

\newcommand{\UserSendsCredit}{El usuario manda crédito}
\newcommand{\UserNoCredit}{El usuario no tiene crédito}
\newcommand{\UserAccountingActivated}{Se activa el accounting para el usuario}

\newcommand{\DataGenerationAll}{El usuario manda, o se generan, todos los datos}
\newcommand{\DataGenerationPartial}{El usuario manda, o se generan, algunos datos}

\newcommand{\ServerlessDown}{La plataforma (serverless) ajusta los recursos a la baja}
\newcommand{\ServerlessUp}{La plataforma (serverless) ajusta los recursos al alza}
\newcommand{\ServerlessAdapts}{La plataforma (serverless) ajusta los recursos según la demanda}

\newcommand{\AccountingIdle}{La plataforma (accounting) no detecta uso así que no consume nada}
\newcommand{\AccountingConsumes}{La plataforma (accounting) detecta uso y va consumiendo el crédito}
\newcommand{\AccountingRestricts}{La plataforma (accounting) restringe los recursos}
\newcommand{\AccountingAllows}{La plataforma (accounting) levanta la restricción}
\newcommand{\AccountingCredit}{El crédito restante debería ser}
\newcommand{\CreditDepleted}{Se agota el crédito}

\newcommand{\AppStart}{Arranca la ejecución de}
\newcommand{\AppRuns}{Continúa la ejecución de la aplicación}
\newcommand{\AppWaits}{La aplicación se para a la espera de que haya crédito}
\newcommand{\AppResume}{La aplicación se reanuda}


\newcommand{\Wait}{Se espera un tiempo}

\newpage
\section{Basic Experiment}

\noindent 
Este experimento busca demostrar de forma específica como se \textbf{aplican y levantan las restricciones} en una misma ejecución y a lo largo del tiempo, la posibilidad de tener una \textbf{deuda}, y de que la \textbf{aplicación no interrumpa su ejecución} en ningún momento (siempre que pueda sobrevivir la restricción).\\

\noindent 
La secuencia sería:
\begin{itemize}
	\setlength\itemsep{0em}
	\setlength{\itemindent}{-2mm}
	\item \DataGenerationAll
	\item \ContainerUp
	\item \UserNoCredit
	\item \UserSendsCredit, poco ($\approx 6~GRC$)
	\item \UserAccountingActivated
	\item \ServerlessDown
	\item \AccountingIdle
	\item \AppStart
	\item \ServerlessUp
	\item \AccountingConsumes
	\item \AppRuns
	\item \UserSendsCredit, una cantidad justa ($ > 4~GRC$), cuando el crédito restante sea bajo ($ < 2~GRC$)
	\item \AppRuns~hasta que se consume todo el crédito 
	\item \CreditDepleted
	\item \AccountingRestricts
	\item \AppRuns~con los recursos restringidos (mínimos)
	\item \AccountingConsumes, se genera una ligera deuda
	\item \UserSendsCredit, una cantidad generosa ($ > 10~GRC$)
	\item \AccountingAllows
	\item \ServerlessAdapts
	\item \AppRuns~hasta su finalización
	\item \ServerlessDown
	\item \AccountingIdle
	\item \AccountingCredit~$ > 0~GRC$
\end{itemize}


\newpage
\section{Trickle processing}


\noindent 
Este experimento busca demostrar de forma específica como el usuario puede manejar la \textbf{ejecución} de su aplicación de forma emparejada con el \textbf{consumo} de crédito, optimizando y controlando dicho gasto en todo momento.\\

\noindent 
La secuencia sería:

\begin{itemize}
	\setlength\itemsep{0em}
	\item \ContainerUp
	\item \UserNoCredit
	\item \UserSendsCredit, muy poco ($\approx 1~GRC$)
	\item \DataGenerationPartial
	\item \UserAccountingActivated
	\item \ServerlessDown
	\item \AccountingIdle
	\item \AppStart
	\item \ServerlessUp
	\item \AccountingConsumes
	\item \AppRuns
	\item De forma periódica y durante la ejecución:
	\begin{itemize}
		\item \DataGenerationPartial
		\item \UserSendsCredit, poco ($ \approx 0.25~GRC$)
	\end{itemize}
	\item El usuario para de mandar crédito y datos
	\item \AppRuns~hasta su finalización
	\item \ServerlessDown
	\item \AccountingIdle
	\item \AccountingCredit~$\approx 0$
\end{itemize}


\newpage
\section{Credit-aware processing}

\noindent 
Este experimento busca demostrar de forma específica como el usuario puede delegar en la aplicación la gestión del consumo de recursos en base al crédito disponible, controlando ahora de manera más indirecta el consumo de recursos y crédito. Hay dos variantes, la de la aplicación que prefiere consumir todo el crédito aunque incurra en una restricción después (greedy), y la que prefiere evitar esto monitorizando el saldo que le queda (polite). Se puede configurar este \textit{umbral de parada} (e.g., 2 GRC)\\

\noindent 
La secuencia sería:

\begin{itemize}
	\setlength\itemsep{0em}
	\item \ContainerUp
	\item \UserNoCredit
	\item \UserAccountingActivated
	\item \AccountingRestricts
	\item \AppStart
	\item \AppWaits~(ahora crédito $= 0~GRC$)
	\item \DataGenerationPartial
	\item \UserSendsCredit~($\approx 10~GRC$)
	\item \AccountingAllows
	\item \AppResume
	\item \ServerlessAdapts
	
	\item \textbf{Escenario Polite}
	\begin{itemize}
		\item \AppRuns, hasta que el crédito esté próximo a agotarse, para evitar una restricción
		\item \AppWaits~(ahora crédito $\approx 2~GRC$)
		\item \ServerlessDown
	\end{itemize}

	\item \textbf{Escenario Greedy}
	\begin{itemize}
		\item \AppRuns~mientras tiene recursos, o sus métricas de rendimiento son buenas, hasta que se consume todo el crédito 
		\item \AccountingRestricts
		\item \AppRuns, pero ahora los recursos son mínimos, la aplicación decide parar
	\end{itemize}

	\item \AccountingIdle
	\item \Wait
	\item \UserSendsCredit~($\approx 10~GRC$)
	\item \DataGenerationPartial

	\item \textbf{Escenario Greedy}
	\begin{itemize}
		\item \AccountingAllows
	\end{itemize}	

	\item \AppResume, ahora que hay crédito disponible
	\item \ServerlessAdapts
	\item \AppRuns, hasta finalizar
	\item \ServerlessDown
	\item \AccountingIdle
	\item \AccountingCredit~$> 0$
\end{itemize}


\newpage

\section{Características de la aplicación}

\begin{itemize}

	\item Poco o nada sensible a cambios de recursos, la aplicación va a tener que seguir ejecutándose aún bajo la restricción de recursos

	\item Con picos y valles de consumo, aunque no debería imponer un alto estrés a la plataforma en su gestión de recursos serverless
		

	\item Sus tareas no pueden tener tiempos de ejecución muy largos ni etapas o fases largas, mejor que tenga una alta granularidad o `capacidad de checkpoining', que se pueda mas o menos parar y reanudar
	
	\item La apliacación leerá los datos de un almacenamiento local (directorio) o externo (e.g., bucket S3, HDFS), pero soportará que dichos datos crezcan, en cuyo caso procesará los nuevos. Los resultados los volcará a otro almacenamiento.
			
\end{itemize}


\section{Candidatos a aplicaciones}


\begin{itemize}
	\item Características comunes:
	\begin{itemize}
		\item Cada candidato se ha probado en un contenedor LXC en mi servidor local
		\item Algunos candidatos usan contenedores Docker, los cuales se arrancan dentro del contenedor LXC
		\item En Plutón no se usa docker sino apptainer, en algunos casos no funcionan las aplicaciones dockerizadas
		\item Para la gestión de datos se puede usar \href{https://min.io/}{minio}, que imita a un servicio S3 en la forma de gestionar buckets y objetos, y se puede desplegar en un contenedor apptainer en local
	\end{itemize}
\end{itemize}

\subsection{Pipeline Metagenómica}

\begin{itemize}
	\item Se usaría el \href{https://gatk.broadinstitute.org/hc/en-us/articles/13832627644443-PathSeqPipelineSpark}{pipeline de GATK para Spark}
	\item Se ejecutaría en local con Spark standalone
	\item Los datos son las muestras de personas y de microorganismos
	\item son 3 etapas
	\begin{itemize}
		\item La primera procesa muestras de microorganismos, hay decenas de muestras, se pueden procesar de forma independiente (de 4 en 4, de 6 en 6), cada muestra consume 1 core al 100\% y tarda desde pocos segundos a +10 minutos. Se usa samtools y Java.
		\item La segunda procesa la muestra humana, hay varias también, se pueden procesar de forma independiente aunque en este caso se usa Spark y se puede especificar el número de cores, por lo que es mejor procesarlas de 1 en 1. Cada muestra tarda, con 32 cores, unos 20 minutos.
		\item La tercera mapea la muestra humana con la de los microorganismos, se usa Spark pero consume menos que la segunda etapa (aprox· 10 minutos) por lo que se podrían hacer varias simultáneamente.
	\end{itemize}
	\item La segunda etapa es intensiva en emoria, requiere $>=$ 64 GiB $\Rightarrow$ Mi servidor local va justo
	\item \textcolor{blue}{\textbf{Se puede usar Plutón}}
	\item \textcolor{blue}{\textbf{Se puede usar mi servidor aunque va algo justo}}
\end{itemize}

\subsection{Suite de benchmarks serverless}

\begin{itemize}
	\item \href{https://github.com/spcl/serverless-benchmarks}{SeBS: Serverless Benchmark Suite}
	\item Se usarían las funciones de compresión (zip), procesamiento de vídeo (ffmpg) y reconocimiento de animales en imagen (pytorch)
	\item Los datos serían carpetas con muchos ficheros a empaquetar y comprimir, vídeos e imágenes de animales
	\item Hay 2 datasets usables que en conjunto suman 11.000 imágenes de animales
	\item La carga viene dada por el número de peticiones y su dificultad, se pueden arrancar varios \textit{workers}, cada uno en un contenedor
	\item Esta suite es la que introdujo minio
	\item \textcolor{red}{\textbf{\textit{Los dockers no los he podido ejecutar con apptainer $\Rightarrow$ No se puede usar Plutón}}}, pero he podido replicar la función de creación de GIF a partir de un MP4 con un contenedor apptainer, se podrían replicar muchas otras (e.g., compresión, encriptado, transcodificación, detección de animales en imagen)
	\item \textcolor{blue}{\textbf{Se puede usar mi servidor local}}
	\item \textcolor{blue}{\textbf{Se puede usar Plutón \textit{con la implementación propia de funciones}}}
\end{itemize}

\subsection{Tienda hecha con microservicios}

\begin{itemize}
	\item La tienda es \href{https://github.com/instana/robot-shop}{Instana Robot Shop}, un modelo de ejemplo de una tienda web con un carrito, productos, pagos
	\item El estrés se genera con \href{https://locust.io/}{Locust}, que simula usuarios haciendo peticiones, clickando en un producto, viendo su descripción, añadiendolo al carrito...
	\item La tienda está formada por una arquitectura de microservicios que se levantan en dockers separados, se ejecutan con Java
	\item En este escenario no se procesan datos, sino peticiones
	\item Locust se puede modular para simular interacción (peticiones) por parte de un número de usuarios
	\item Puede ser necesario arrancar varios estresadores Locust porque sino se ven limitados por CPU
	\item Algún microservicio consume bastante memoria (+10 GiB)
	\item \textcolor{blue}{\textbf{Se puede usar mi servidor local}}
	\item \textcolor{red}{\textbf{La ejecución de los docker con apptainer no me ha funcionado por algo relacionado con usar el usuario root así como no usar el entrypoint $\Rightarrow$ No se puede usar Plutón (de momento)}}
\end{itemize}




\subsection{Procesado de eventos tipo clickstream}

\begin{itemize}
	\item Se usaría \href{https://github.com/ClickHouse/ClickBench}{ClickBench}
	\item El dataset se carga en una MongoDB (otras DDBB soportadas) y se emiten queries
	\item Las queries tienen duraciones distintas, las hay instantáneas, de 1, 10, 20, 30 minutos
	\item Las queries cargan los workers de la base de datos que consumen CPU
	\item Las queries también son intensivas en disco 
	\item \textcolor{blue}{\textbf{Se puede usar mi servidor local}}
	\item \textcolor{black}{\textbf{Se puede usar Plutón? (pendiente de comprobar)}}
	\item \textcolor{red}{\textbf{Las queries son intensivas en disco, se requiere mínimo un SSD que soporte 250 MBps (medición empírica, un SSD de gama media/alta debería llegar) para que no tarden demasiado}}
\end{itemize}


\subsection{Procesado de series temporales}

\begin{itemize}
	\item Se usaría la suite \href{https://github.com/timescale/tsbs}{Time Series Benchmark Suite (TSBS)}
	\item Los datos se generan con la propia herramienta, hay 2 escenarios, devops e iot
	\item El escenario devops simula servidores y métricas
	\item El escenario iot simula camiones en ruta y eventos
	\item La base de datos usada seria PostgreSQL con la extensión para series temporales TimescaleDB
	\item La carga es modulable con una variable `scale', en devops el número de hosts, en iot el de camiones
	\item La carga puede ser la propia inserción de datos, o la ejecución de queries
	\item Las queries se generan también con la herramienta y pueden modularse de igual manera con una variable `scale'
	\item \textcolor{red}{\textbf{No he podido meter PostgreSQL en un contenedor apptainer y sospecho que no se pueda hacer de forma directa dado que postgresql requiere ejecutarse con el usuario postgres y escribir sus datos en un directorio del usuario postgres $\Rightarrow$ No se puede usar Plutón}}
	\item \textcolor{blue}{\textbf{Se puede usar mi servidor local}}
	\item \textcolor{red}{\textbf{Experimentos intensos en disco}}
\end{itemize}


\subsection{Otros}
\begin{itemize}
		
	\item Reconocimiento facial (e.g., \href{http://vis-www.cs.umass.edu/lfw/}{Faces in the Wild}), en este caso los datos son las caras.
		
	\item Procesado de ficheros de texto (e.g., \href{https://github.com/logpai}{logs, LogPai})
	
\end{itemize}


\end{document}
